% !TEX encoding = UTF-8 Unicode
% !TEX program  = xelatex
\documentclass[12pt,aspectratio=169]{beamer}

% ====== FONTS ======
\usepackage{fontspec}

% ====== THEME (Metropolis: flat, modern, minimalist) ======
\usetheme[numbering=fraction,progressbar=frametitle]{metropolis}

% Re-assert fonts AFTER the theme
\setmainfont{Arial}
\setsansfont{Arial}
\setmonofont{Consolas}[Scale=0.88]

\usepackage{tikz}
\usetikzlibrary{positioning,arrows.meta,shapes.geometric,fit,calc,shadows.blur,decorations.pathreplacing}
\usepackage{booktabs}
\usepackage{array}
\usepackage{tabularx}
\usepackage[table]{xcolor}
\usepackage{graphicx}

% ====== COLOR PALETTE ======
\definecolor{ecInk}{HTML}{1B2430}      % Near-black text
\definecolor{ecPrimary}{HTML}{1E3A8A}  % Deep indigo blue
\definecolor{ecPrimaryDk}{HTML}{142a63}
\definecolor{ecAccent}{HTML}{F2784B}   % Warm coral accent
\definecolor{ecTeal}{HTML}{12877F}     % Soft teal secondary
\definecolor{ecPaper}{HTML}{FAFAF8}    % Warm off-white background
\definecolor{ecPanel}{HTML}{EEF1F8}    % Light indigo panel fill
\definecolor{ecPanelAlt}{HTML}{E4F3F1} % Light teal panel fill
\definecolor{ecLine}{HTML}{D7DCE6}     % Hairline / table rule

\setbeamercolor{normal text}{fg=ecInk,bg=ecPaper}
\setbeamercolor{alerted text}{fg=ecAccent}
\setbeamercolor{example text}{fg=ecTeal}
\setbeamercolor{palette primary}{fg=ecPaper,bg=ecPrimary}
\setbeamercolor{frametitle}{fg=ecPaper,bg=ecPrimary}
\setbeamercolor{progress bar}{fg=ecAccent,bg=ecLine}
\setbeamercolor{progress bar in head/foot}{fg=ecAccent,bg=ecLine}
\setbeamercolor{progress bar in section page}{fg=ecAccent,bg=ecLine}
\setbeamercolor{title separator}{fg=ecAccent,bg=ecAccent}
\setbeamercolor{block title}{fg=ecPrimary,bg=ecPanel}
\setbeamercolor{block body}{fg=ecInk,bg=ecPanel!55!white}
\setbeamercolor{section title}{fg=ecPrimary}
\setbeamercolor{footnote}{fg=ecInk!70}
\setbeamercolor{caption}{fg=ecInk!70}

\setbeamerfont{frametitle}{size=\large,series=\bfseries}
\setbeamerfont{title}{size=\LARGE,series=\bfseries}
\setbeamerfont{subtitle}{size=\normalsize,series=\mdseries}
\setbeamerfont{block title}{size=\normalsize,series=\bfseries}

\setbeamertemplate{frametitle continuation}{}
\makeatletter
\setlength{\metropolis@titleseparator@linewidth}{2pt}
\setlength{\metropolis@progressonsectionpage@linewidth}{2pt}
\setlength{\metropolis@progressinheadfoot@linewidth}{1.4pt}
\makeatother

% ====== TABLE STYLE ======
\renewcommand{\arraystretch}{1.25}
\newcolumntype{L}[1]{>{\raggedright\arraybackslash}p{#1}}
\newcolumntype{Y}{>{\raggedright\arraybackslash}X}

% Reusable components
\newcommand{\kw}[2]{{\color{#1}\textbf{#2}}}

\newcommand{\statpill}[2]{%
  \tikz[baseline=-0.6ex]{
    \node[draw=ecAccent,line width=1pt,rounded corners=6pt,fill=ecAccent!12,
          inner sep=4pt,font=\bfseries\scriptsize,text=ecPrimaryDk] {#1 #2};
  }%
}
\newcommand{\slideimage}[3]{%
  \includegraphics[width=#2\textwidth,height=#3\textheight,keepaspectratio]{#1}%
}

\title{ĐẶC TẢ HỆ THỐNG \textbf{EmotiCare AIoT}}
\subtitle{Thiết bị AIoT thông minh đồng hành và chăm sóc sức khỏe cảm xúc}
\author{Trần Hải Đức \,\textbar\, Nguyễn Trọng Tài \,\textbar\, Nguyễn Công Chiến \,\textbar\, Phạm Nguyễn Gia Bảo}
\institute{Trường Đại học Khoa học Tự nhiên --- ĐHQG TP.HCM\\ Khoa Công nghệ Thông tin \,\textbar\, Lớp 23CLC02}
\date{Specification v3.2 \,\textbar\, Tháng 6/2026}

\begin{document}

% ===================== SLIDE TRANG BÌA =====================
\begin{frame}[plain]
  \begin{tikzpicture}[remember picture,overlay]
    \fill[ecPrimary] (current page.south west) rectangle (current page.north east);
    \fill[ecAccent] ([yshift=-1.05cm]current page.north west) rectangle ([yshift=-1.13cm]current page.north east);
    \node[anchor=north east,xshift=-0.9cm,yshift=-0.55cm,text=ecPaper!65,font=\scriptsize\bfseries\sffamily]
      at (current page.north east) {EMOTICARE \textbar\ AIoT SPEC};
  \end{tikzpicture}
  \vspace{1.3cm}
  {\setbeamercolor{titlelike}{fg=ecPaper}%
   \setbeamercolor{title}{fg=ecPaper}%
   \setbeamercolor{subtitle}{fg=ecPaper!85}%
   \setbeamercolor{author}{fg=ecPaper!92}%
   \setbeamercolor{date}{fg=ecPaper!75}%
   \setbeamercolor{institute}{fg=ecPaper!75}%
   \titlepage}
\end{frame}

% ===================== MỤC LỤC =====================
\begin{frame}{Nội dung trình bày}
  \tableofcontents
\end{frame}

% ===================== SECTION 00: GIỚI THIỆU =====================
\section{Giới thiệu sản phẩm}

\begin{frame}{Định vị sản phẩm}
  \begin{block}{EmotiCare AIoT là gì?}
    Thiết bị AIoT giúp người dùng \textbf{nhận biết, thấu hiểu và quản lý cảm xúc} trong cuộc sống hằng ngày thông qua giọng nói.
  \end{block}
  \vspace{0.35cm}
  \kw{ecPrimary}{Cách tiếp cận:} Edge-first nhưng Cloud-assisted
  \begin{itemize}
    \item \textbf{Edge AI} nhận diện cảm xúc cốt lõi tại thiết bị để phản hồi nhanh và hạn chế lộ dữ liệu giọng nói.
    \item \textbf{Cloud} hỗ trợ các tính năng cần Internet như gợi ý hoạt động, nhạc, podcast, trò chuyện và báo cáo dài hạn.
    \item \textbf{TFT screen} là điểm theo dõi chính: hiển thị cảm xúc, gợi ý, trạng thái sync và báo cáo ngày/tháng/năm.
  \end{itemize}
  \vfill
  \centering {\color{ecPrimary}\small\itshape ``EmotiCare AIoT -- Understand your feelings, care for your mind.''}
\end{frame}

\begin{frame}{Mục đích sản phẩm}
  \small
  \begin{block}{Tóm tắt mục đích}
    Giúp người dùng \textbf{nhận biết cảm xúc hiện tại}, \textbf{nhận hỗ trợ đúng lúc} và \textbf{theo dõi xu hướng cảm xúc ngay trên TFT}.
  \end{block}
  \vspace{0.1cm}
  \begin{itemize}
    \item Nhận diện cảm xúc trong \statpill{15s}{} sau một tương tác giọng nói hợp lệ
    \item Đề xuất hoạt động, bài hát, podcast hoặc phản hồi đồng cảm trong \statpill{20s}{}
    \item Lưu emotion session để phục vụ đồng bộ, phân tích và xem lại lịch sử
    \item Tạo báo cáo cảm xúc theo \textbf{ngày, tháng, năm} và trả về TFT trong \statpill{180s}{}
  \end{itemize}
\end{frame}

% ===================== SECTION 01: BACKGROUND =====================
\section{Background}

\begin{frame}{Bối cảnh \& Vấn đề}
  \begin{block}{Thực trạng}
    \small Nhu cầu chăm sóc sức khỏe tinh thần cao, nhưng người dùng thường khó ghi nhận cảm xúc đều đặn; self-care cần một điểm chạm đơn giản, riêng tư và dùng được hằng ngày.
  \end{block}
  \vspace{0.05cm}
  {\renewcommand{\arraystretch}{0.96}
  \scriptsize
  \rowcolors{2}{ecPanel!45}{white}
  \begin{tabularx}{\textwidth}{L{2.9cm}YY}
  \toprule
  \rowcolor{ecPrimary}\textcolor{white}{\textbf{Vấn đề}} & \textcolor{white}{\textbf{Hệ quả}} & \textcolor{white}{\textbf{Giải pháp EmotiCare}} \\
  \midrule
  Ít ghi nhận cảm xúc & Khó thấy xu hướng dài hạn & Check-in giọng nói, lưu emotion session \\
  Cảm xúc tiêu cực kéo dài & Dễ bỏ lỡ dấu hiệu căng thẳng & Báo cáo ngày/tháng/năm trên TFT \\
  Gợi ý chăm sóc chung chung & Khó chọn hoạt động phù hợp & Cá nhân hóa theo cảm xúc và feedback \\
  Dữ liệu giọng nói nhạy cảm & Người dùng lo ngại riêng tư & Edge AI, không upload audio thô mặc định \\
  \bottomrule
  \end{tabularx}}
  \vspace{0.05cm}
  \centering\scriptsize\color{ecInk!65}\textit{Nguồn tham khảo trong specification: WHO World Mental Health Report, NIMH self-care guidance.}
\end{frame}

\begin{frame}{Nguồn cảm hứng \& Sản phẩm tương tự}
  \small
  \begin{columns}[T,onlytextwidth]
    \column{0.48\textwidth}
    \kw{ecPrimary}{Nguồn cảm hứng:} EMO (LivingAI)
    \begin{itemize}
      \item AI desktop pet tạo cảm giác đồng hành gần gũi
      \item EmotiCare kế thừa tinh thần đó, nhưng chuyển trọng tâm sang \textbf{chăm sóc cảm xúc có dữ liệu}
    \end{itemize}
    \begin{center}
      \slideimage{images/emo.jpg}{0.40}{0.20}
    \end{center}
    \column{0.48\textwidth}
    \kw{ecTeal}{So sánh nhanh}
    \begin{itemize}
      \item \textbf{EMO}: giải trí, cá tính, chưa tập trung SER
      \item \textbf{ElliQ}: companion/wellness cho người lớn tuổi
      \item \textbf{EmotiCare}: Edge SER + Cloud support + báo cáo TFT
    \end{itemize}
  \end{columns}
  \vfill
  \centering\color{ecAccent}\textbf{Định vị:} hỗ trợ tự chăm sóc cảm xúc, không chẩn đoán hoặc thay thế chuyên gia.
\end{frame}

\begin{frame}{Người dùng mục tiêu}
  \small
  \rowcolors{2}{ecPanel!45}{white}
  \begin{tabularx}{\textwidth}{L{3.5cm}YY}
  \toprule
  \rowcolor{ecPrimary}\textcolor{white}{\textbf{Nhóm}} & \textcolor{white}{\textbf{Nhu cầu chính}} & \textcolor{white}{\textbf{Giá trị sản phẩm}} \\
  \midrule
  Sinh viên & Theo dõi căng thẳng, áp lực học tập & Check-in nhanh, gợi ý nghỉ ngơi \\
  Người đi làm & Quản lý stress trong ngày làm việc & Nhận biết thời điểm căng thẳng \\
  Người quan tâm wellness & Xây thói quen chăm sóc tinh thần & Báo cáo định kỳ \\
  Gia đình/người chăm sóc & Góc nhìn tổng quan (nếu được chia sẻ) & Báo cáo không xâm phạm riêng tư \\
  \bottomrule
  \end{tabularx}
\end{frame}

\begin{frame}{Từ mục đích suy ra 3 SMART Objective}
  \begin{center}
    \slideimage{images/Goal To Objective.png}{0.90}{0.75}
  \end{center}
\end{frame}

% ===================== SECTION 02: KIẾN TRÚC =====================
\section{Kiến trúc \& Phần cứng}

\begin{frame}{Kiến trúc tổng quan: Edge -- Cloud -- TFT}
  \begin{center}
    \slideimage{images/Architechture.png}{0.80}{0.55}
  \end{center}
  \vspace{-0.1cm}
  \begin{itemize}
    \item \kw{ecPrimary}{Edge}: thu âm, chạy SER, hiển thị kết quả, hoạt động offline
    \item \kw{ecTeal}{Cloud}: gợi ý, hội thoại, báo cáo --- cần Internet
    \item \kw{ecAccent}{TFT}: giao diện theo dõi duy nhất của người dùng
  \end{itemize}
\end{frame}

\begin{frame}{Thành phần chính của hệ thống}
  \small
  {\renewcommand{\arraystretch}{1.1}
  \rowcolors{2}{ecPanel!45}{white}
  \begin{tabularx}{\textwidth}{L{3.8cm}Y}
  \toprule
  \rowcolor{ecPrimary}\textcolor{white}{\textbf{Thành phần}} & \textcolor{white}{\textbf{Vai trò}} \\
  \midrule
  Edge Device & Mic, nút bấm, Wi-Fi, cache, TFT \\
  Edge AI SER Engine & Nhận diện cảm xúc từ giọng nói (cục bộ) \\
  Cloud API Server & Cổng giao tiếp Edge $\leftrightarrow$ Cloud \\
  Recommendation Service & Gợi ý hoạt động, bài hát, podcast \\
  Conversation Service & Phản hồi đồng cảm, có safety filter \\
  Report Engine & Tổng hợp báo cáo cảm xúc theo thời gian \\
  Cloud Database & Lưu trữ dài hạn (user, session, report,...) \\
  \bottomrule
  \end{tabularx}}
\end{frame}

\begin{frame}{Phần cứng đề xuất \& chi phí prototype}
  \small
  \rowcolors{2}{ecPanel!45}{white}
  \begin{tabularx}{\textwidth}{Y L{3.5cm}}
  \toprule
  \rowcolor{ecPrimary}\textcolor{white}{\textbf{Linh kiện}} & \textcolor{white}{\textbf{Giá tham khảo}} \\
  \midrule
  ESP32-S3 (bộ điều khiển chính) & 180k -- 350k VND \\
  INMP441 I2S Microphone & 25k -- 70k VND \\
  TFT/OLED Display & 120k -- 320k VND \\
  5 nút bấm/touch + Speaker & 20k -- 110k VND \\
  \rowcolor{ecAccent!18}\textbf{Tổng dự đoán (đầy đủ)} & \textbf{545k -- 1.700k VND} \\
  \bottomrule
  \end{tabularx}
  \vspace{0.2cm}
  \begin{block}{Nguyên tắc thiết kế}
    Edge cho nhận diện cốt lõi --- Cloud cho hỗ trợ nâng cao --- TFT là giao diện chính --- Chịu lỗi offline --- An toàn cảm xúc.
  \end{block}
\end{frame}

% ===================== SECTION 03: OBJECTIVES =====================
\section{3 SMART Objectives \& Use Case}

\begin{frame}{Tổng quan 3 SMART Objective}
  \small
  \rowcolors{2}{ecPanel!45}{white}
  \begin{tabularx}{\textwidth}{L{2.8cm}Y L{3.2cm}}
  \toprule
  \rowcolor{ecPrimary}\textcolor{white}{\textbf{Objective}} & \textcolor{white}{\textbf{Mô tả đầy đủ}} & \textcolor{white}{\textbf{Vai trò trong vòng lặp}} \\
  \midrule
  SMART Objective 1 & Phát hiện và phân loại cảm xúc trong 15s sau tương tác giọng nói hợp lệ; lưu emotion session & Tạo dữ liệu nền cho hỗ trợ và báo cáo \\
  SMART Objective 2 & Đề xuất hoạt động, bài hát, podcast hoặc phản hồi đồng cảm trong 20s khi có Internet & Biến cảm xúc/nhu cầu từ HOME thành hành động hỗ trợ \\
  SMART Objective 3 & Tạo tóm tắt thống kê theo ngày/tháng/năm trên Cloud và trả report cards về TFT trong 180s & Giúp người dùng nhìn lại xu hướng dài hạn \\
  \bottomrule
  \end{tabularx}
\end{frame}

\begin{frame}{UC-01: Speech Emotion Recognition}
  \textbf{Input:} Giọng nói \quad $\rightarrow$ \quad \textbf{Output:} Trạng thái cảm xúc (vui vẻ, bình thường, căng thẳng, buồn bã, tức giận, mệt mỏi)
  \vspace{0.25cm}
  \begin{enumerate}
    \item Người dùng kích hoạt Check-in
    \item Edge AI tiền xử lý âm thanh, trích xuất đặc trưng
    \item Mô hình SER phân loại cảm xúc + confidence
    \item TFT hiển thị kết quả, lưu emotion session
  \end{enumerate}
  \begin{block}{Mục tiêu hiệu năng}
    Hoàn tất trong \statpill{15s}{}. Hoạt động được cả khi \textbf{offline}.
  \end{block}
\end{frame}

\begin{frame}{UC-02, UC-03, UC-04: Hỗ trợ qua Cloud (20s)}
  \small
  \rowcolors{2}{ecPanel!45}{white}
  \begin{tabularx}{\textwidth}{L{1.5cm}L{3.5cm}Y}
  \toprule
  \rowcolor{ecPrimary}\textcolor{white}{\textbf{UC}} & \textcolor{white}{\textbf{Tên}} & \textcolor{white}{\textbf{Mô tả ngắn}} \\
  \midrule
  UC-02 & Gợi ý hoạt động \& nội dung & Cloud đề xuất hoạt động, bài hát, podcast theo emotion label/lịch sử \\
  UC-03 & Chọn bài hát/podcast theo chủ đích & Người dùng chọn category (thư giãn, tập trung, ngủ nghỉ,...) \\
  UC-04 & Trò chuyện hỗ trợ cảm xúc & Phản hồi đồng cảm, có safety filter, xử lý tín hiệu nguy cấp \\
  \bottomrule
  \end{tabularx}
  \vspace{0.2cm}
  \centering\small\color{ecInk!70}\textit{Tất cả đều cần Internet; nếu mất kết nối, TFT báo trạng thái cần Cloud.}
\end{frame}

\begin{frame}{UC-05: Thống kê \& phân tích xu hướng cảm xúc}
  \textbf{Input:} Emotion sessions, activity logs, media logs, conversation metadata \\
  \textbf{Output:} Báo cáo rút gọn theo ngày/tháng/năm trên TFT
  \vspace{0.25cm}
  \begin{itemize}
    \item Cloud Report Engine tính phân bố cảm xúc, xu hướng, hiệu quả hoạt động/nội dung
    \item Nén thành các \textbf{TFT report cards} ngắn gọn
    \item Nếu dữ liệu ít $\rightarrow$ hiển thị trạng thái \texttt{limited\_data}
  \end{itemize}
  \begin{block}{Mục tiêu hiệu năng}
    Hiển thị trên TFT trong vòng \statpill{180s}{}.
  \end{block}
\end{frame}

\begin{frame}{Overall Use Case Diagram}
  \begin{center}
    \slideimage{images/overall UC.png}{0.88}{0.72}
  \end{center}
\end{frame}

% ===================== SECTION 04: EDGE AI =====================
\section{Edge AI -- Speech Emotion Recognition}

\begin{frame}{Pipeline SER trên Edge Device}
  \begin{center}
    \slideimage{images/EdgeAI Pipeline.png}{0.90}{0.40}
  \end{center}
  \vspace{0.15cm}
  \begin{itemize}
    \item Dữ liệu tham khảo: \textbf{RAVDESS} (Kaggle) -- tập giọng nói cảm xúc có nhãn
    \item Đặc trưng chính: Log-Mel Spectrogram, MFCC, pitch, energy
    \item Mô hình: CNN nhỏ (có thể kèm LSTM/GRU nhẹ), tối ưu bằng quantization/TFLite Micro
  \end{itemize}
\end{frame}

\begin{frame}{Ánh xạ nhãn cảm xúc (RAVDESS $\rightarrow$ sản phẩm)}
  \small
  \rowcolors{2}{ecPanel!45}{white}
  \begin{tabularx}{\textwidth}{L{3.5cm}L{3.0cm}Y}
  \toprule
  \rowcolor{ecPrimary}\textcolor{white}{\textbf{Nhãn RAVDESS}} & \textcolor{white}{\textbf{Nhãn sản phẩm}} & \textcolor{white}{\textbf{Ý nghĩa}} \\
  \midrule
  neutral, calm & Bình thường & Trạng thái ổn định \\
  happy & Vui vẻ & Cảm xúc tích cực \\
  sad & Buồn bã & Cần phản hồi đồng cảm \\
  angry & Tức giận & Cần gợi ý tạm dừng, thở chậm \\
  fearful, surprised & Căng thẳng & Cần hỗ trợ giảm áp lực \\
  -- & Mệt mỏi & Nhãn mở rộng, cần dữ liệu bổ sung \\
  \bottomrule
  \end{tabularx}
\end{frame}

\begin{frame}{Logic confidence \& quality flag}
  \small
  \rowcolors{2}{ecPanel!45}{white}
  \begin{tabularx}{\textwidth}{L{6.0cm}Y}
  \toprule
  \rowcolor{ecPrimary}\textcolor{white}{\textbf{Điều kiện}} & \textcolor{white}{\textbf{Hành vi hệ thống}} \\
  \midrule
  confidence $\ge$ 0.75 \& clean & Hiển thị emotion label, chuyển sang hỗ trợ \\
  0.50 $\le$ confidence $<$ 0.75 & Hiển thị ``có thể là...'', cho xác nhận \\
  confidence $<$ 0.50 & Không chắc chắn, không kết luận xu hướng \\
  quality\_flag = too\_short & Yêu cầu ghi âm lại \\
  quality\_flag = noisy & Cảnh báo môi trường nhiễu \\
  \bottomrule
  \end{tabularx}
  \vspace{0.15cm}
  \centering\scriptsize\color{ecInk!70}\textit{Đánh giá mô hình: Accuracy, Confusion matrix, Macro F1, Latency, Model size, Robustness test.}
\end{frame}

% ===================== SECTION 05: INTERNET SERVICE =====================
\section{Internet Service}

\begin{frame}{Sequence: Edge -- Cloud -- TFT}
  \begin{center}
    \slideimage{images/Internet 2.png}{0.88}{0.72}
  \end{center}
\end{frame}

\begin{frame}{Thiết kế Database (rút gọn)}
  \small
  \textbf{Bảng chính:} \texttt{users}, \texttt{devices}, \texttt{emotion\_sessions}, \texttt{recommendation\_requests}, \texttt{activity\_feedback}, \texttt{media\_items}, \texttt{media\_selection\_logs}, \texttt{conversation\_requests}, \texttt{tft\_reports}
  \vspace{0.3cm}
  
  {\renewcommand{\arraystretch}{1.2}
  \rowcolors{2}{ecPanel!45}{white}
  \begin{tabularx}{\textwidth}{L{3.8cm}Y}
  \toprule
  \rowcolor{ecPrimary}\textcolor{white}{\textbf{Tên bảng}} & \textcolor{white}{\textbf{Thuộc tính tiêu biểu}} \\
  \midrule
  \texttt{emotion\_sessions} & \texttt{client\_session\_id (UNIQUE), emotion\_label, confidence\_score, quality\_flag, sync\_status} \\
  \midrule
  \rowcolor{ecTeal!15}\texttt{tft\_reports} & \texttt{period\_type, tft\_cards (JSONB), emotion\_distribution, data\_quality} \\
  \bottomrule
  \end{tabularx}}
\end{frame}

\begin{frame}{API tiêu biểu cho Edge Device}
  \small
  \rowcolors{2}{ecPanel!45}{white}
  \begin{tabularx}{\textwidth}{Y L{1.8cm} Y}
  \toprule
  \rowcolor{ecPrimary}\textcolor{white}{\textbf{Endpoint}} & \textcolor{white}{\textbf{Method}} & \textcolor{white}{\textbf{Mục đích}} \\
  \midrule
  \texttt{/api/emotion-sessions/sync} & POST & Đồng bộ emotion sessions (idempotent) \\
  \texttt{/api/recommendations/request} & POST & Lấy activity/song/podcast cards \\
  \texttt{/api/conversations/respond} & POST & Lấy response card \\
  \texttt{/api/reports/tft-summary} & GET & Lấy report cards theo period \\
  \bottomrule
  \end{tabularx}
\end{frame}

% ===================== SECTION 06: FR/NFR =====================
\section{Functional \& Non-Functional Requirements}

\begin{frame}{Tóm tắt Functional Requirements}
  \small
  \rowcolors{2}{ecPanel!45}{white}
  \begin{tabularx}{\textwidth}{L{3.2cm}Y L{2.0cm}}
  \toprule
  \rowcolor{ecPrimary}\textcolor{white}{\textbf{Nhóm}} & \textcolor{white}{\textbf{Nội dung chính}} & \textcolor{white}{\textbf{ID}} \\
  \midrule
  Nhận diện Edge & Check-in, ghi âm, SER, lưu session trong 15s & FR-01--08 \\
  Đồng bộ nền tảng & Cache offline, sync idempotent, trạng thái TFT & FR-09--13 \\
  Gợi ý qua Cloud & Activity cards, song cards, podcast cards trong 20s & FR-14--27 \\
  Trò chuyện qua Cloud & Phản hồi đồng cảm, safety filter & FR-28--34 \\
  Báo cáo qua Cloud & Report cards theo ngày/tháng/năm trong 180s & FR-35--41 \\
  \bottomrule
  \end{tabularx}
\end{frame}

\begin{frame}{Tóm tắt Non-Functional Requirements}
  \small
  \rowcolors{2}{ecPanel!45}{white}
  \begin{tabularx}{\textwidth}{L{3.5cm}Y}
  \toprule
  \rowcolor{ecPrimary}\textcolor{white}{\textbf{Nhóm}} & \textcolor{white}{\textbf{Yêu cầu chính}} \\
  \midrule
  Hiệu năng & SER $\le$ 15s, gợi ý/hội thoại $\le$ 20s, report $\le$ 180s \\
  Tin cậy \& khả dụng & Offline cho Objective 1, retry sync, idempotency \\
  Bảo mật \& riêng tư & Không upload audio thô mặc định, device token, HTTPS \\
  An toàn cảm xúc & Không chẩn đoán, phản hồi đồng cảm, xử lý tín hiệu nguy cấp \\
  Trải nghiệm TFT & Thao tác đơn giản, report cards ngắn gọn dễ đọc \\
  \bottomrule
  \end{tabularx}
\end{frame}

% ===================== SECTION 07: USER MANUAL =====================
\section{User Manual -- Screen Flow}

\begin{frame}{Luồng màn hình thiết bị}
  \scriptsize
  \rowcolors{2}{ecPanel!45}{white}
  \renewcommand{\arraystretch}{1.05}
  \begin{tabularx}{\textwidth}{L{3.0cm}Y}
  \toprule
  \rowcolor{ecPrimary}\textcolor{white}{\textbf{Màn hình}} & \textcolor{white}{\textbf{Luồng thao tác chính}} \\
  \midrule
  HOME & Chọn Check-in, Activity, Music/Podcast, Conversation, Report hoặc Status \\
  CHECK-IN $\rightarrow$ RESULT & Thu giọng nói, chạy SER tại Edge, hiển thị emotion label và confidence \\
  SUPPORT & Sau check-in, chọn hoạt động, nội dung nghe hoặc trò chuyện \\
  ACTIVITY / MEDIA / CHAT & Có Internet để lấy activity cards, media cards hoặc response card từ Cloud \\
  REPORT & Chọn ngày/tháng/năm; xem report từ Cloud hoặc cache gần nhất \\
  STATUS & Xem online/offline, pending sessions và lần sync gần nhất \\
  \bottomrule
  \end{tabularx}
\end{frame}

\begin{frame}{Ví dụ trải nghiệm sử dụng}
  \begin{enumerate}
    \item Từ \textbf{HOME}, người dùng có thể chọn thẳng Check-in, Activity, Music/Podcast, Conversation hoặc Report.
    \item Nếu chọn \textbf{Check-in}, người dùng nói một câu ngắn vào thiết bị.
    \item TFT hiển thị ngay: \textit{Căng thẳng -- confidence 0.74}.
    \item Chuyển nhanh sang \textbf{Activity}: Gợi ý ``Hít thở 4-7-8'' + đề xuất nghe podcast thở chậm.
    \item Người dùng thực hiện và đánh giá hoạt động $\rightarrow$ phản hồi gửi trực tiếp lên Cloud.
    \item Mở màn hình \textbf{Report}: Chọn xem theo ngày/tháng/năm để nắm bắt tỷ lệ phân bố cảm xúc và hoạt động hữu ích nhất.
  \end{enumerate}
\end{frame}

% ===================== SECTION 08: KẾT LUẬN =====================
\section{Kết luận}

\begin{frame}{Tổng kết \& Vòng lặp vận hành}
  \centering
  \vskip 0.5cm
  {\large\color{ecPrimary}\bfseries Check-in $\rightarrow$ Edge SER $\rightarrow$ Hiển thị TFT $\rightarrow$ Đồng bộ Cloud $\rightarrow$ Hỗ trợ/Báo cáo $\rightarrow$ Hiển thị TFT}
  \vspace{0.5cm}
  \begin{itemize}
    \item \textbf{Edge AI} đảm nhiệm xử lý nhận diện cảm xúc cốt lõi, hoạt động tin cậy bất kể trạng thái mạng.
    \item \textbf{Cloud} mở rộng tính năng với các dịch vụ gợi ý thông minh, hội thoại sâu và phân tích báo cáo dài hạn.
    \item \textbf{TFT} giữ vai trò là điểm chạm phần cứng duy nhất, trực quan và bảo mật quyền riêng tư tối đa.
  \end{itemize}
\end{frame}

\begin{frame}{Lợi ích \& Giới hạn}
  \begin{columns}[T,onlytextwidth]
    \column{0.48\textwidth}
    \begin{block}{Lợi ích}
      \begin{itemize}\small
        \item Tăng khả năng tự nhận thức cảm xúc.
        \item Hỗ trợ đúng lúc, theo dõi dài hạn.
        \item Phù hợp tối ưu cho mô hình prototype sinh viên.
        \item Riêng tư cao (không upload âm thanh gốc).
      \end{itemize}
    \end{block}
    \column{0.48\textwidth}
    \begin{alertblock}{Giới hạn}
      \begin{itemize}\small
        \item SER mang tính xác suất, có tỉ lệ sai số.
        \item Objective 2, 3 phụ thuộc hoàn toàn vào Internet.
        \item Không gian hiển thị trên TFT khá hạn chế.
        \item Không có chức năng thay thế thiết bị y tế.
      \end{itemize}
    \end{alertblock}
  \end{columns}
\end{frame}

\begin{frame}{Hướng phát triển tiếp theo}
  \begin{itemize}
    \item Học mô hình baseline cảm xúc cá nhân hóa riêng cho từng đối tượng người dùng.
    \item Liên tục cập nhật và nén các mô hình SER tối ưu hơn chuyên biệt cho các dòng Edge Device.
    \item Nâng cao thuật toán cá nhân hóa gợi ý dựa trên lịch sử tương tác và hiệu quả phản hồi.
    \item Tối ưu hóa ngôn ngữ thiết kế đồ họa nhỏ gọn, trực quan trên màn hình TFT (mini-charts).
    \item Tích hợp bản đồ tài nguyên và liên hệ hỗ trợ tâm lý khẩn cấp theo định vị khu vực.
  \end{itemize}
\end{frame}

% ===================== SLIDE KẾT THÚC =====================
\begin{frame}[plain]
  \begin{tikzpicture}[remember picture,overlay]
    \fill[ecPrimary] (current page.south west) rectangle (current page.north east);
  \end{tikzpicture}
  \centering
  \vfill
  {\color{white}\Large\bfseries Cảm ơn quý thầy/cô và các bạn đã theo dõi!}
  \vskip 0.4cm
  {\color{ecAccent}\large\texttt{EmotiCare AIoT -- Intelligent Emotional Companion}}
  \vskip 1.5cm
  {\color{white!80!black}\small Trần Hải Đức \,\textbar\, Nguyễn Trọng Tài \,\textbar\, Nguyễn Công Chiến \,\textbar\, Phạm Nguyễn Gia Bảo}
  \vfill
\end{frame}

\end{document}
